\section{Introducción: de lo discreto a lo continuo}

En las clases anteriores, hemos estudiado los fundamentos de los modelos de difusión (Diffusion Models), incluyendo el proceso hacia adelante y el proceso hacia atrás del DDPM (Denoising Diffusion Probabilistic Models), así como el método del DDIM (Denoising Diffusion Implicit Models) para convertir un proceso estocástico en uno determinista con el fin de reducir el número de pasos de muestreo. Estos conocimientos básicos son suficientes para construir y entrenar sus propios modelos de difusión en muchas aplicaciones prácticas.

El objetivo central de esta clase es comprender los modelos de difusión desde una nueva perspectiva: el \textbf{dominio continuo del tiempo}. Veremos que:

\begin{itemize}
    \item Tanto el DDPM de tiempo discreto como el SMLD (Score Matching con Dinámica de Langevin) pueden unificarse bajo un marco de Ecuaciones Diferenciales Estocásticas (SDE, Stochastic Differential Equation);
    \item Dada una SDE hacia adelante, se puede derivar directamente la SDE hacia atrás;
    \item Existe una forma equivalente a la SDE en términos de Ecuaciones Diferenciales Ordinarias (ODE, Ordinary Differential Equation), denominada ODE de flujo de probabilidad;
    \item Este marco continuo nos otorga mayor libertad y herramientas teóricas para diseñar modelos de difusión.
\end{itemize}

\begin{knowledgebox}{Antecedentes del curso}
Esta clase se basa en el artículo clásico de Song et al. (2021) titulado ``Score-Based Generative Modeling through Stochastic Differential Equations''. Este trabajo unifica los dos paradigmas SMLD (NCSN) y DDPM, proponiendo un marco general de SDE. Es una literatura fundamental para comprender la teoría moderna de los modelos de difusión.
\end{knowledgebox}

\begin{importantbox}{Puntos clave}
Anteriormente definimos el proceso hacia adelante y hacia atrás de los modelos de difusión bajo \textbf{tiempo discreto}. El cambio clave en esta clase es que, cuando el número de pasos discretos $N \to \infty$, todo el proceso de adición de ruido hacia adelante puede describirse mediante una SDE, mientras que el proceso de eliminación de ruido hacia atrás corresponde igualmente a una SDE inversa. Esta perspectiva de tiempo continuo no solo unifica el SMLD y el DDPM, sino que también proporciona la base teórica para diseñar nuevos modelos de difusión y muestreadores eficientes.
\end{importantbox}
